\documentclass[11pt]{amsart}
\usepackage{amsmath,amssymb,amsthm}
\usepackage[margin=1.15in]{geometry}
\usepackage[hidelinks]{hyperref}
\hypersetup{pdftitle={Major roads and span length in scale-invariant random spatial networks},pdfauthor={Yitzchak Shmalo}}
\newtheorem{theorem}{Theorem}[section]
\newtheorem{lemma}[theorem]{Lemma}
\newtheorem{corollary}[theorem]{Corollary}
\theoremstyle{remark}
\newtheorem{remark}[theorem]{Remark}
\newcommand{\R}{\mathbb{R}}
\newcommand{\Z}{\mathbb{Z}}
\newcommand{\E}{\mathbb{E}}
\newcommand{\Prob}{\mathbb{P}}
\newcommand{\len}{\operatorname{len}}
\newcommand{\spn}{\operatorname{span}}
\newcommand{\doi}[1]{\href{https://doi.org/#1}{doi:\nolinkurl{#1}}}
\newcommand{\arxiv}[1]{\href{https://arxiv.org/abs/#1}{arXiv:\nolinkurl{#1}}}
\setlength{\emergencystretch}{2.5em}

\title[Major roads in a scale-invariant random spatial network]{Major roads and span length in scale-invariant
random spatial networks}
\author{Yitzchak Shmalo}
\address{Einstein Institute of Mathematics, The Hebrew University of Jerusalem, Jerusalem, Israel}
\email{yspennstate@gmail.com}
\date{November 2026}
\subjclass[2020]{60D05, 60G55, 90B15}
\keywords{scale-invariant random spatial network, SIRSN, binary hierarchy model, geodesics, Poisson process}

\raggedbottom
\begin{document}
\begin{abstract}
We study two of Aldous's open problems on scale-invariant random spatial networks. For his binary hierarchy model we
show that almost surely the set of major roads contains every road line above a fixed height. These lines form a
connected grid, so the major roads have exactly one unbounded connected component, which answers Open Problem~7 of
Aldous for this model. Open Problem~9 asks under what extra assumptions the subnetwork spanning $k$ uniform points in a
square of area $k$ has expected length asymptotic to $\ell k$. We show, in every such network, that this holds exactly
when the part of the span far from the square has expected length $o(k)$, and that it holds when the expected length
of one route outside the disc of radius $t$ about an endpoint is $o(t^{-3})$. The condition is met in the binary
hierarchy model, with no extra assumption, and in the planar Poisson line network when its speed exponent exceeds~$5$.
\end{abstract}
\maketitle

\section{Introduction}

A scale-invariant random spatial network (SIRSN), in the sense of Aldous~\cite{A14}, assigns to every pair of points
$z,z'$ of the plane a route $\mathcal R(z,z')$, a path of finite length from $z$ to $z'$, so that the collection of
routes is invariant in law under translation, rotation and scaling, routes between the points of a Poisson process
use a network of finite edge-intensity, and long routes do not wander. Aldous and Ganesan~\cite{AG13} gave the
motivation from road networks, and Aldous~\cite{A14} built the first example, the binary hierarchy model, by placing
roads of height $h$ along the lines of $2^h\Z$ that are not in $2^{h+1}\Z$, with travel cost $\gamma^{h}$ per unit
length, $1/2<\gamma<1$, and randomising the origin, the orientation and the scale. Kendall~\cite{K17} showed that a
Poisson process of lines with power-law speed limits gives a random metric space and a pre-SIRSN in the plane, and
Kahn~\cite{Kahn16} proved that it is a SIRSN in every dimension. Aldous's paper ends with a list of open problems; this
paper treats two of them. We use the numbering of the published version~\cite{A14}. In the arXiv version the two
questions are Open Problems~33 and~35, and the results of~\cite{A14} cited below as Lemmas~3.2, 3.3, 3.7 and~3.10 and
Propositions~3.1 and~3.8 are Lemmas~3, 4, 8 and~11 and Propositions~2 and~9.

Let $\Pi$ be a Poisson process of intensity $dz\,dt$ on $\R^2\times[0,\infty)$, independent of the network, and let
$\Xi(\lambda)$ be the set of $z$ with $(z,t)\in\Pi$ for some $t\le\lambda$. Then $\Xi(\lambda)$ is a Poisson process of
rate $\lambda$, increasing in $\lambda$. Let $\mathcal E(\lambda,r)$ be the union, over $\xi,\xi'\in\Xi(\lambda)$, of
$\mathcal R(\xi,\xi')\setminus(\bar B(\xi,r)\cup\bar B(\xi',r))$, where $\bar B(z,r)$ is the closed disc of radius
$r$~\cite[(2.17)]{A14}. Aldous calls $\mathcal E(\infty,1)=\bigcup_\lambda\mathcal E(\lambda,1)$ the major roads. He
notes that $\mathcal E(\infty,1)$ need not be connected but must contain an unbounded component, and asks whether it
has only one~\cite[Open Problem~7]{A14}. Theorem~\ref{thm:roads} answers this for the binary hierarchy model.

The proof uses Aldous's lemmas on routes in that model~\cite[Lemmas~3.2, 3.3 and~3.7]{A14} and his continuum
construction~\cite[Proposition~3.8]{A14}: the route crosses the highest vertical line in its range in one segment,
stays in a box of comparable size, and costs at most a power of the distance. The proof of Lemma~3.10 of~\cite{A14},
which Aldous credits to Justin Salez for its details, already bounds the excess cost of a path below by a multiple of
the length of a road segment it misses. Lemma~\ref{lem:straddle} adapts this comparison, using cheap connecting routes
at the two ends and the confinement of routes, to a quantitative statement uniform in the lattice scale: a route
between points close to a high road line, on opposite sides of it, runs along that line for almost its whole height.
So every point of every high line lies on a route between points far from it, and the high lines form a grid.

For the Poisson road model, Blanc, Curien and Kahn~\cite[Theorem~2 and Propositions~4 and~5]{BCK25} identify the
geodesic frame with the union of the road lines. In their statement the distance from a road point to the endpoints of
the geodesic through it may depend on the point, whereas $\mathcal E(\infty,1)$ requires the fixed distance~$1$, so
their theorem does not settle Open Problem~7 for that model. Other properties of the Poisson road metric are studied
by Blanc~\cite{Bl24} and Kendall~\cite{K20}.

Aldous also asks~\cite[Open Problem~9]{A14} under what extra assumptions, if any, the expected total length of the
subnetwork spanned by $k$ uniform points in a square of area $k$ is asymptotic to $\ell k$, where $\ell$ is the
edge-intensity of the subnetwork on a rate-one Poisson process. Theorem~\ref{thm:span} combines Aldous's bound on the
intensity of $\mathcal E(\lambda,r)$ with a comparison of the uniform sample to a Poisson process and a dyadic
decomposition of the plane outside the square. This controls the span at every distance up to $D$ from the square at
cost $O(\sqrt k\log D+D)$ and reduces what is left to a tail estimate for one route. The estimate holds automatically
in the binary hierarchy model by Aldous's bounded stretch~\cite[Proposition~3.1]{A14}, and in the planar Poisson line
model with speed exponent above $5$ by Kahn's moment bounds~\cite[proof of Theorem~5.1]{Kahn16}
(Corollary~\ref{cor:span}).

\section{The major roads of the binary hierarchy model}\label{sec:roads}

We use the model and notation of~\cite[Section~3]{A14}. For an integer $H$ let $\Z^2_H$ denote the lattice of
spacing $2^H$. On it a coordinate $c\in2^H\Z$ has height $h$ if $c\in2^h\Z\setminus2^{h+1}\Z$, and $0$ has height $\infty$. A lattice path
has cost equal to the sum over its segments of length times $\gamma^{h}$, where $h$ is the height of the line carrying
the segment and $\gamma^\infty=0$. Routes are paths of minimal cost; among paths of equal cost the route is chosen by
Aldous's secondary randomisation~\cite[Section~3.4]{A14}, which makes the choices consistent. We use four facts
from~\cite[Section~3]{A14}. The first three hold for every path of minimal cost and, by the binary scaling of the
model, on every $\Z^2_H$ with the same constants.
\begin{itemize}
\item[(F1)] If $I_x$ is the range of $x$-coordinates of a route and $X$ is the unique point of maximal height in
$I_x$, the route meets the line $\{x=X\}$ in a single segment or a single point~\cite[Lemma~3.2(ii)]{A14}.
\item[(F2)] The route from $z_1$ to $z_2$ stays inside the square of side $K\|z_2-z_1\|_1$ centred at $z_1$, where
$K=K(\gamma)$~\cite[Lemma~3.3]{A14}.
\item[(F3)] The minimal cost from $z_1$ to $z_2$ is at most $K'\|z_2-z_1\|_1^\beta$, where $2^\beta=2\gamma$, so
$0<\beta<1$~\cite[Lemma~3.7]{A14}. On $\Z^2_H$ the cost of a path is $(2\gamma)^{H}$ times the cost on $\Z^2$ of the
path scaled by $2^{-H}$, so $K'$ does not depend on $H$.
\item[(F4)] For $z_1,z_2$ outside a Lebesgue-null set $A$ there is a continuum route $\mathcal R_0(z_1,z_2)$, the limit
of the lattice routes between lattice approximations $z_{iH}\to z_i$ as $H\to-\infty$: for each $H_0$ the part of
the lattice route lying on lines of height at least $H_0$ is the same for all sufficiently small $H$, and it is
contained in $\mathcal R_0(z_1,z_2)$~\cite[Proposition~3.8]{A14}.
\end{itemize}

\begin{lemma}\label{lem:straddle}
Let $X$ have height $h$ and let $z_1=(x_1,y_1)$, $z_2=(x_2,y_2)$ be points of $\Z^2_H$ with $x_1<X<x_2$,
$|x_i-X|\le\delta$ and $y_2-y_1=2a>0$. Suppose $K(a+\delta)+\delta<2^h$ and
\[
\eta:=\frac{2K'\delta^\beta}{\gamma^{h-1}(1-\gamma)}<a .
\]
Then every route from $z_1$ to $z_2$ contains the segment $\{X\}\times[y_1+\eta,\,y_2-\eta]$. The same holds if $X$
has height larger than $h$.
\end{lemma}

\begin{proof}
By (F2) the range $I_x$ lies within distance $K(a+\delta)+\delta<2^h$ of $X$. The points of height at least $h$ form
$2^h\Z$, so $X$ is the only such point in $I_x$, and it is the peak. By (F1) the route meets $\{x=X\}$ in one segment
(or point). Let $e$ denote the height at which the route enters it and $f$ the height at which it leaves, in the order
of traversal from $z_1$. Before reaching
$(X,e)$ the route changes its height by at least $|e-y_1|$ on vertical lines other than $X$, and after leaving $(X,f)$
by at least $|y_2-f|$; every vertical line in $I_x$ other than $X$ has height at most $h-1$. Let
$S=|e-y_1|+|y_2-f|$. Then
\[
\operatorname{cost}\ \ge\ |f-e|\,\gamma^{h}+S\,\gamma^{h-1}.
\]
Concatenate a route from $z_1$ to $(X,y_1)$, the segment from $(X,y_1)$ to $(X,y_2)$, and a route from $(X,y_2)$ to
$z_2$. By (F3) this path costs at most $2K'\delta^\beta+2a\gamma^h$, and deleting any loops does not increase its cost.
Since $|f-e|\ge 2a-S$, minimality gives $S(\gamma^{h-1}-\gamma^h)\le2K'\delta^\beta$, that is $S\le\eta$. If $e>f$
then $S\ge 2a+(e-f)>\eta$, which is impossible; so $e\le f$, $e\le y_1+\eta$ and $f\ge y_2-\eta$, and the route
contains the stated segment, which is nonempty because $\eta<a$. If $X$ has height $h_X>h$, replace $\gamma^h$ by
$\gamma^{h_X}$ on the line $X$; since $\gamma^{h-1}-\gamma^{h_X}\ge\gamma^{h-1}(1-\gamma)$, the same bound holds.
\end{proof}

\begin{theorem}\label{thm:roads}
In the binary hierarchy model let $h_0$ be the least integer with $2^{h_0}>6K+1$. Almost surely $\mathcal E(\infty,1)$
contains every road line of height at least $h_0$, where heights are the labels the lines inherit from the lattice
construction (defined in the proof). These lines form a connected square grid whose spacing is at most $2^{h_0+1}$.
Every connected component of $\mathcal E(\infty,1)$ of diameter larger than $\sqrt2$ times that spacing is the
component containing the grid, and $\mathcal E(\infty,1)$ has exactly one unbounded connected component.
\end{theorem}

\begin{proof}
\emph{One line of the lattice model.} For $h\ge h_0$ let
\[
\delta_h=\min\Bigl(\frac18,\Bigl[\frac{\gamma^{h-1}(1-\gamma)}{16K'}\Bigr]^{1/\beta}\Bigr),
\]
so that the quantity $\eta$ of Lemma~\ref{lem:straddle} with $\delta=\delta_h$ satisfies $\eta\le1/8$. Let $X$ have
height $h\ge h_0$, let $y\in\R$, and let $\xi,\xi'\notin A$ lie in the open discs of radius $\delta_h/8$ about
$(X-\delta_h/2,\,y-4)$ and $(X+\delta_h/2,\,y+4)$. For all sufficiently small $H$ the lattice approximations of
$\xi,\xi'$ lie within $\delta_h/8$ of them, so their $x$-coordinates lie in $(X-\delta_h,X)$ and $(X,X+\delta_h)$, and
half their vertical separation $a_H$ satisfies $|a_H-4|\le\delta_h/4$. Then
$K(a_H+\delta_h)+\delta_h<6K+1<2^h$ and $\eta\le1/8<a_H$, so by Lemma~\ref{lem:straddle} every lattice route between
the approximations contains $\{X\}\times[y_{1}+\eta,\,y_{2}-\eta]$, where $y_1\le y-4+\delta_h/4$ and
$y_2\ge y+4-\delta_h/4$ are the second coordinates of the approximations. This segment contains
$\{X\}\times[y-1,y+1]$ and lies on a line of height $h$, so by (F4) the continuum route $\mathcal R_0(\xi,\xi')$
contains $\{X\}\times[y-1,y+1]$. Each point of this segment is at distance more than $2$ from $\xi$ and from $\xi'$.
Horizontal lines are treated in the same way.

\emph{The translation.} Following~\cite[Section~3.7]{A14}, let $U_n$ be uniform on $[0,2^n]^2$, coupled so that
$U_n=U_{n+1}\bmod 2^n$ coordinatewise, and let
$\mathcal R^{(n)}(z_1,z_2)=\mathcal R_0(z_1+U_n,z_2+U_n)-U_n$. For fixed $z_1,z_2$, almost surely the translation
invariant network satisfies $\mathcal R_{ti}(z_1,z_2)=\mathcal R^{(n)}(z_1,z_2)$ for all sufficiently large
$n$~\cite[(3.3), (3.4)]{A14}. For $h\in\Z$ and $n>h$ put $B^x_h=2^h\Z-U^x_n$ and $B^y_h=2^h\Z-U^y_n$; since
$U_m-U_n\in2^n\Z^2$ for $m\ge n$, these sets do not depend on $n$. They decrease in $h$, and a coordinate in
$B^x_h\setminus B^x_{h+1}$ is said to have label $h$. For $n>h+1$ the translation by $-U_n$ carries the vertical lines
of height $h$ of the lattice model onto the vertical lines with label $h$. Almost surely no coordinate lies in every
$B^x_h$: a bounded interval $I$ meets $B^x_h$ with probability at most $|I|2^{-h}$, these events decrease in $h$, and
countably many intervals cover $\R$. So the vertical lines of label at least $h_0$ are the lines $\{x=c\}$ with
$c\in B^x_{h_0}$, a translate of $2^{h_0}\Z$, and the same holds for horizontal lines.

\emph{The endpoints.} The network $\mathcal R$ is obtained from $\mathcal R_{ti}$ by a uniform rotation and a scaling
by a factor in $(1/2,2]$~\cite[Section~3.7]{A14}, applied to points and routes together and independent of $(U_n)$.
Let $D'$ be the image of $\bigcup_\lambda\Xi(\lambda)$ under the inverse of these two maps. Conditional on the
rotation and the scale, $D'$ is an increasing union of locally finite Poisson samples whose intensities tend to
infinity, so it is almost surely countable and dense. The space-time sampling process is independent of the hierarchy,
including its secondary randomisation, so conditional on these and on the rotation and the scale, $D'$ almost surely
misses the null set $\bigcup_n(A-U_n)$. Conditional on $D'$, the stabilisation of the previous paragraph
holds for every pair of points of $D'$ outside an event of probability zero, since there are countably many pairs.
On the intersection of these events let $\{x=c\}$ have label $h\ge h_0$, let $y$ be rational, and choose
$\xi,\xi'\in D'$ in the open discs of radius $\delta_h/8$ about $(c-\delta_h/2,\,y-4)$ and $(c+\delta_h/2,\,y+4)$.
For $n>h+1$ the coordinate $c+U^x_n$ has height $h$ in the lattice model, and the first step applied to
$\xi+U_n$, $\xi'+U_n$ shows that $\mathcal R^{(n)}(\xi,\xi')$ contains $\{c\}\times[y-1,y+1]$. For large $n$ this
route is $\mathcal R_{ti}(\xi,\xi')$. As $y$ runs over the rationals these segments cover the line, and there are
countably many lines of label at least $h_0$. After the rotation and the scaling, each point of these segments lies on
the route between two points of $\bigcup_\lambda\Xi(\lambda)$ at distance more than $2\cdot\frac12=1$ from both, hence
in $\mathcal E(\infty,1)$. Horizontal lines are treated in the same way.

\emph{The grid.} Let $G$ be the image of the union of the lines of label at least $h_0$. It is a connected and
unbounded square grid of spacing $s\le2^{h_0+1}$, contained in $\mathcal E(\infty,1)$. The complement of $G$ is a union
of disjoint open squares of side $s$, so a connected subset of $\mathcal E(\infty,1)$ that misses $G$ lies inside one
of them and has diameter at most $\sqrt2\,s$. Every component of larger diameter, in particular every unbounded
component, meets $G$ and therefore contains it, and the component containing $G$ is unbounded.
\end{proof}

\begin{remark}
The proof shows more than Open Problem~7 asks: the major roads contain a connected set $G$ whose complement has only
bounded components. For a general SIRSN this property would give uniqueness without any statement on the coalescence
of semi-infinite geodesics. In the Poisson line model of Kendall~\cite{K17} and Kahn~\cite{Kahn16} the natural
candidate for $G$ is the union of the lines above some speed, and what is needed is that every point of such a line
lies on a route between points at distance more than one from it. The theorem of Blanc, Curien and
Kahn~\cite{BCK25} gives each point of a road line on a geodesic, with endpoints at a distance that depends on the point.
\end{remark}

\section{The length of the spanning subnetwork}\label{sec:span}

Now let $\mathcal R$ be any SIRSN. We use the notation and axioms of~\cite{A14}: $\mathcal S(\lambda)$ is the union of
the routes between points of $\Xi(\lambda)$, with edge-intensity $\ell\lambda^{1/2}$, where
$\ell=\E\len(\mathcal S(1)\cap[0,1]^2)<\infty$; $p(\lambda,r)$ is the edge-intensity of $\mathcal E(\lambda,r)$ and
$p(r)=\lim_{\lambda\to\infty}p(\lambda,r)$; the axioms give $p(1)<\infty$ and $D_1=\len\mathcal R(0,(1,0))$ with
$\E D_1<\infty$. By scaling, $p(\lambda,r)=p(\lambda r^2,1)/r$, and since $\mathcal E(\lambda,r)$ increases with
$\lambda$ we get $p(\lambda,r)\le p(r)=p(1)/r$. Define the one-route tail
\[
h(t)=\E\,\len\bigl(\mathcal R(0,(1,0))\setminus B(0,t)\bigr),\qquad t\ge0 ,
\]
where $B(0,t)$ is the open disc and $B(0,0)=\emptyset$. It decreases to $0$ by dominated convergence. Let $Q$ be a
square of side $L=\sqrt k$, let $Z_1,\dots,Z_k$ be independent and uniform on $Q$, independent of the network, and let
$\spn(Z_1,\dots,Z_k)$ be the union of the routes $\mathcal R(Z_i,Z_j)$. In this section the constants implied by
$O(\cdot)$ and $o(\cdot)$ depend only on the law of the network, and limits are taken as $k\to\infty$.

\begin{theorem}\label{thm:span}
For every $D=D(k)\ge1$,
\[
\E\len\spn(Z_1,\dots,Z_k)\ \le\ \ell k+O(k^{2/3})+p(1)\bigl(4L\log_2D+6\pi D\bigr)+O(L)
+\tfrac{1}{\sqrt2}\,k^2L\,h\!\left(\tfrac{D}{\sqrt2L}\right),
\]
and $\E\len\bigl(\spn(Z_1,\dots,Z_k)\cap Q\bigr)\ge(\ell-o(1))k$. Consequently:
\begin{enumerate}
\item $\E\len\spn(Z_1,\dots,Z_k)=\ell k+o(k)$ if and only if, for some $D(k)\ge1$ with $D(k)=o(k)$, the expected length
of the part of the span at distance at least $D(k)$ from $Q$ is $o(k)$;
\item if $h(t)=o(t^{-3})$ as $t\to\infty$, then $\E\len\spn(Z_1,\dots,Z_k)\sim\ell k$.
\end{enumerate}
\end{theorem}

\begin{proof}
\emph{Upper bound.} Let $\lambda=1+k^{-1/3}$, let $U_1,U_2,\dots$ be independent and uniform on $Q$, independent of
$N\sim\mathrm{Poisson}(\lambda k)$, and put $Z_i=U_i$. Then $\{U_1,\dots,U_N\}$ has the law of $\Xi(\lambda)\cap Q$,
and we couple the two. On $\{N\ge k\}$ the span of $Z_1,\dots,Z_k$ lies in $\mathcal S(\lambda)$. Split the span into
its parts inside $Q$, at distance in $(0,1]$ from $Q$, at distance in $(2^j,2^{j+1}]$ for $0\le j\le J=\lfloor\log_2
D\rfloor$, and at distance at least $D$.

Inside $Q$, the expected length of $\mathcal S(\lambda)\cap Q$ is $\ell\lambda^{1/2}k=\ell k+O(k^{2/3})$. Within
distance $1$ of $Q$ it is at most $\ell\lambda^{1/2}$ times the area $4L+\pi$ of that ring, which is $O(L)$. A point of
the span at distance $d\in(2^j,2^{j+1}]$ from $Q$ lies on a route between two points of $Q$ at distance more than $2^j$
from both, hence in $\mathcal E(\lambda,2^j)$, whose intensity is at most $p(1)2^{-j}$. The set of points at distance
at most $r$ from $Q$ has area $L^2+4Lr+\pi r^2$, so the shell has area $4L2^j+3\pi4^j$ and its expected length is at
most $p(1)(4L+3\pi2^{j})$. Summing over $0\le j\le J$ gives
$p(1)\bigl(4L(J+1)+3\pi(2^{J+1}-1)\bigr)\le p(1)(4L\log_2D+6\pi D)+4p(1)L$.

On $\{N<k\}$ we use the trivial bound: the length of the union of all the portions within distance $D$ of $Q$ is at
most $\len\spn(Z_1,\dots,Z_k)$, which is independent of $N$ and has mean at most $\sum_{i<j}\E\len\mathcal R(Z_i,Z_j)\le k^2\sqrt2L\,\E D_1$ by scaling. Since
$\Prob(N<k)\le\exp(-ck^{1/3})$, this contributes $o(1)$.

Beyond distance $D$: a point of $\mathcal R(Z_i,Z_j)$ at distance at least $D$ from $Q$ is at distance at least $D$
from $Z_i$. By scaling and rotation invariance,
$\E\len(\mathcal R(z,z')\setminus B(z,D))=|z-z'|\,h(D/|z-z'|)$, and $s\mapsto s\,h(D/s)$ is nondecreasing in $s$, so for
$z,z'\in Q$ this is at most $\sqrt2L\,h(D/(\sqrt2L))$. A union bound over the at most $k^2/2$ pairs gives the last
term.

\emph{Lower bound inside $Q$.} Fix $R\ge1$ and $M>0$, and retain the unit squares $c$ of a tiling of $Q$ whose
$R$-neighbourhood lies in $Q$; there are $k-O(R\sqrt k)$ of them. In such a square the span contains every route
between sample points in the $R$-neighbourhood $N_R(c)$ of $c$. Conditional on their number, these points are
independent and uniform on $N_R(c)$, and their number is binomial with mean the area $A_R$ of $N_R(c)$, which is within
total variation $O(A_R^2/k)$ of a Poisson variable. Since the network is independent of the sample, the expectation of
the minimum of $M$ and the length in $c$ of those routes is within $O(MA_R^2/k)$ of the same quantity for $\Xi(1)$,
uniformly over the retained squares by stationarity. Let $k\to\infty$, then $R\to\infty$, then $M\to\infty$. Each point
of $\mathcal S(1)$ lies on a route between a finite pair of points, so the quantity for $\Xi(1)$ increases with $R$ to
$\E\min(M,\len(\mathcal S(1)\cap c))$, and then with $M$ to $\ell$. All these lengths lie inside $Q$, so
$\liminf_k k^{-1}\E\len(\spn\cap Q)\ge\ell$.

\emph{Consequences.} Let $T_k=\len\spn(Z_1,\dots,Z_k)$ and $I_k=\len(\spn(Z_1,\dots,Z_k)\cap Q)$. The upper bound
gives $\E I_k\le\ell k+o(k)$ and the lower bound $\E I_k\ge\ell k-o(k)$. If $\E T_k=\ell k+o(k)$, then
$\E(T_k-I_k)=o(k)$, so the condition in (1) holds with $D=1$. Conversely, if the part of the span at distance at least
$D(k)$ from $Q$ has expected length $o(k)$ for some $D(k)\ge1$ with $D(k)=o(k)$, the bounds for the parts within
distance $D$ give $\E T_k\le\ell k+O(k^{2/3})+O(\sqrt k\log D+D)+o(k)=\ell k+o(k)$, and $\E T_k\ge\E I_k\ge\ell k-o(k)$.
This proves (1). For (2), let $\varepsilon(t)=\sup_{u\ge t}u^3h(u)$, which is finite for large $t$ and tends to $0$,
and for large $k$ put
\[
\eta_k=\max\bigl(k^{-1/4},\,\varepsilon(k^{1/4}/\sqrt2)^{1/6}\bigr),\qquad D=\eta_kk,\qquad t_k=\frac{D}{\sqrt2L}
=\frac{\eta_k\sqrt k}{\sqrt2}.
\]
Then $\eta_k\to0$, so $1\le D=o(k)$, and $t_k\ge k^{1/4}/\sqrt2$. The last term of the upper bound divided by $k$ is
$\tfrac1{\sqrt2}k^{3/2}h(t_k)\le\tfrac1{\sqrt2}k^{3/2}\varepsilon(t_k)t_k^{-3}=2\eta_k^{-3}\varepsilon(t_k)$, which is at
most $2\varepsilon(k^{1/4}/\sqrt2)^{1/2}\to0$ (and is $0$ if $\varepsilon(k^{1/4}/\sqrt2)=0$). Together with the
lower bound this gives (2).
\end{proof}

\begin{corollary}\label{cor:span}
\begin{enumerate}
\item In the binary hierarchy model, $(\ell-o(1))k\le\E\len\spn(Z_1,\dots,Z_k)\le\ell k+O(k^{2/3})$.
\item If $\E D_1^4<\infty$, then $\E\len\spn(Z_1,\dots,Z_k)\sim\ell k$. This holds for the planar Poisson line SIRSN
of Kendall~\cite{K17} and Kahn~\cite{Kahn16} when its speed exponent, called $\gamma$ in~\cite{Kahn16}, exceeds~$5$.
\end{enumerate}
\end{corollary}

\begin{proof}
A route from $0$ that leaves $B(0,t)$ has length at least $t$, so $h(t)\le\E(D_1\mathbf 1_{\{D_1\ge t\}})$, where
$\mathbf 1_{\{D_1\ge t\}}$ is the indicator of the event $\{D_1\ge t\}$. In the binary hierarchy model, Aldous's bounded
stretch~\cite[Proposition~3.1]{A14}, with its constant $K_\gamma$, gives
$\len\mathcal R_0(z_1,z_2)\le\sqrt2K_\gamma|z_2-z_1|$, a bound unaffected by the transformations taking $\mathcal R_0$
to $\mathcal R$~\cite[Section~3.7, p.~21]{A14}; so $D_1\le\sqrt2K_\gamma$ and $h(t)=0$ for $t>\sqrt2K_\gamma$. Taking
$D=3K_\gamma L$ in Theorem~\ref{thm:span} makes the last term vanish and the middle terms $O(\sqrt k\log k)$, which
proves (1). If $\E D_1^4<\infty$, then $t^3h(t)\le\E(D_1^4\mathbf 1_{\{D_1\ge t\}})\to0$, and (2) follows from
Theorem~\ref{thm:span}(2). For the Poisson line model, Kahn's proof of~\cite[Theorem~5.1]{Kahn16} shows that the
Euclidean length of the geodesic between two fixed points has finite moments of every order below $\gamma-1$.
\end{proof}

\begin{remark}
The axioms of~\cite{A14} control the union of all routes through the intensity bound $p(1)/r$; they do not supply an
explicit rate for $h$ in their formulation, and the argument above uses an additional tail estimate. The factor $k^2$
in the last term comes from the union bound over pairs, and far pieces of different routes largely coincide, so a
count of distinct far route segments would weaken the tail condition. For the Poisson line SIRSN with speed exponent
between $2$ and $5$ the question is how often a fastest path between nearby points uses a distant fast line.
\end{remark}

\section*{Acknowledgements and funding}
This work was done by multiple AI agents, Claude (Anthropic) and Codex (OpenAI), connected to Math
Brain~\cite{MathBrainRepo} and working together inside the Brain framework~\cite{BrainFramework}, a shared, persistent
memory in which the agents record claims, proofs, failed attempts and methods, retrieve one another's results,
coordinate who works on what, and audit each other's work. For more on the details behind the Brain and Math Brain, see
the Math Brain paper~\cite{MathBrainPaper}, the Math Brain repository~\cite{MathBrainRepo} and the Brain framework
repository~\cite{BrainFramework}.

The research presented in this paper was supported by the European Research Council (ERC) under the European Union's
Horizon Europe research and innovation programme (grant agreement No.~101041711), by the Simons Foundation as part of
the Collaboration on the Mathematical and Scientific Foundations of Deep Learning, by Heights Labs, by Convex Nexus
Capital, by the Israel Science Foundation (grant number 2258/19), and by the Israel Science Foundation (ISF Grant
4101/25). Computational resources were provided in part by Google Cloud through its education credits programme (credit
530550103).

\section*{Data availability}
No result of the paper depends on computation. The \LaTeX{} source of the paper is available at\\
\url{https://github.com/yspennstate/sirsn-major-roads-binary-hierarchy/tree/v1.0}.

The paper is part of the Math Brain project, whose public repository is given in~\cite{MathBrainRepo}.

\section*{Declaration of competing interest}
The author holds equity in Heights Labs and in Convex Nexus Capital and has a management role at Convex Nexus Capital.
These entities are included in the funding declaration above.

\section*{Declaration of generative AI and AI-assisted technologies}
Claude agents (Anthropic) proved Theorems~\ref{thm:roads} and~\ref{thm:span} and prepared the manuscript; a Codex agent
(OpenAI) refereed the proofs independently. The SIRSN axioms, the binary hierarchy model, the route lemmas quoted as
(F1)--(F4), the bounded stretch estimate and the two questions are due to David Aldous~\cite{A14}, with motivation from
Aldous and Karthik Ganesan~\cite{AG13}; Aldous credits Justin Salez for the details of the proof of his Lemma~3.10,
whose excess-cost comparison precedes the straddling argument here. Wilfrid S. Kendall~\cite{K17} and Jonas
Kahn~\cite{Kahn16} established the Poisson road construction and its SIRSN property, including the moment bounds used
in Corollary~\ref{cor:span}, and Guillaume Blanc, Nicolas Curien and Kahn~\cite{BCK25} proved the related
geodesic-frame theorem. New here are the straddling bound of Lemma~\ref{lem:straddle}, its fixed-distance grid
consequence for the binary hierarchy model in Theorem~\ref{thm:roads}, and the span bound of Theorem~\ref{thm:span}
with Corollary~\ref{cor:span}. The author directed the work and is responsible for the content of the paper and for
decisions concerning its dissemination.

\begin{thebibliography}{99}
\bibitem{A14} D.~J. Aldous, \emph{Scale-invariant random spatial networks}, Electron. J. Probab. \textbf{19} (2014),
no.~15, 1--41, \doi{10.1214/EJP.v19-2920}; arXiv version \arxiv{1204.0817}, 2012, with different numbering.
\bibitem{AG13} D.~J. Aldous and K.~Ganesan, \emph{True scale-invariant random spatial networks}, Proc. Natl. Acad. Sci.
USA \textbf{110} (2013), 8782--8785, \doi{10.1073/pnas.1304329110}.
\bibitem{Bl24} G.~Blanc, \emph{Fractal properties of Aldous--Kendall random metric}, Ann. Inst. Henri Poincar\'e
Probab. Stat. \textbf{60} (2024), 1356--1386, \doi{10.1214/22-AIHP1354}, \arxiv{2207.03349}.
\bibitem{BCK25} G.~Blanc, N.~Curien and J.~Kahn, \emph{Geodesics in planar Poisson roads random metric}, Proc. Lond.
Math. Soc. \textbf{131} (2025), no.~1, e70070, \doi{10.1112/plms.70070}, \arxiv{2407.07887}.
\bibitem{Kahn16} J.~Kahn, \emph{Improper Poisson line process as SIRSN in any dimension}, Ann. Probab. \textbf{44}
(2016), 2694--2725, \doi{10.1214/15-AOP1032}, \arxiv{1503.03976}.
\bibitem{K17} W.~S. Kendall, \emph{From random lines to metric spaces}, Ann. Probab. \textbf{45} (2017), 469--517,
\doi{10.1214/14-AOP935}, \arxiv{1403.1156}.
\bibitem{K20} W.~S. Kendall, \emph{Rayleigh random flights on the Poisson line SIRSN}, Electron. J. Probab.
\textbf{25} (2020), no.~124, 1--36, \doi{10.1214/20-EJP526}.
\bibitem{MathBrainRepo} Y.~Shmalo, \emph{Math Brain}, public software repository, snapshot of 14 November 2026,
\url{https://github.com/yspennstate/math-brain-public/tree/8c9a6bd478c4ec4082347151b7ae16db50c917dd}.
\bibitem{MathBrainPaper} Y.~Shmalo, \emph{Math Brain: shared, auditable memory for multi-agent mathematical research},
manuscript dated 26 October 2026, Zenodo release 12 November 2026,
\url{https://github.com/yspennstate/math-brain-public/tree/main/papers/math-brain-2026-10-24},
\doi{10.5281/zenodo.24150837}.
\bibitem{BrainFramework} Y.~Shmalo, \emph{Brain Framework}, public software repository, snapshot of 13 November 2026,
\url{https://github.com/yspennstate/brain-framework-public/tree/5629e644cdd510173ae97e9aa92642dabd22dc57}.
\end{thebibliography}

\end{document}
